% ==================================================
% Difference Rule — Consequences and Examples
% Discrete Calculus — Polynomial Basis
% ==================================================

\section*{Consequences of the Difference Rule}

The difference rule \(\Delta\,x^{\underline{k}} = k\,x^{\underline{k-1}}\)
(Theorem~\ref{thm:falling-factorial-rule}) has several important consequences.

\section*{Higher Differences of Falling Factorials}

Applying the rule repeatedly gives:
\[
\Delta^m\, x^{\underline{k}}
=
\begin{cases}
\dfrac{k!}{(k-m)!}\, x^{\underline{k-m}} & \text{if } m \le k, \\[6pt]
0 & \text{if } m > k.
\end{cases}
\]
In particular, evaluating at \(x = 0\):
\[
\Delta^m\, x^{\underline{k}}\big|_{x=0}
=
\begin{cases}
k! & \text{if } m = k, \\
0 & \text{if } m \ne k.
\end{cases}
\]

\begin{remark*}[Basis orthogonality]
This vanishing property means the falling factorials behave like a
``triangular'' basis: \(\Delta^m\) kills all \(x^{\underline{k}}\) with
\(k < m\), and returns a scalar multiple for \(k = m\). This is what
makes the coefficient extraction in Newton's expansion work cleanly.
\end{remark*}

\section*{Connection to Binomial Coefficients}

The higher difference formula from the Binomial Theorem section
(Proposition~\ref{prop:higher-diff-formula}),
\[
\Delta^n f(x) = \sum_{k=0}^{n}(-1)^{n-k}\binom{n}{k}f(x+k),
\]
together with the difference rule, gives an explicit formula for the
coefficients of the falling factorial expansion:
\[
c_k = \frac{\Delta^k f(0)}{k!}.
\]

\begin{remark*}[Pascal's triangle revisited]
The coefficients appearing in higher finite differences follow exactly
the pattern of Pascal's triangle — a direct consequence of the
algebraic identity \(\Delta = E - I\) and the binomial theorem.
\end{remark*}

\section*{Worked Example: Expanding \(n^2\)}

We find the falling factorial expansion of \(f(n) = n^2\).

Computing the difference table at \(n = 0\):
\[
f(0) = 0, \quad
\Delta f(0) = 1, \quad
\Delta^2 f(0) = 2, \quad
\Delta^3 f(0) = 0.
\]

The coefficients are \(c_k = \Delta^k f(0)/k!\):
\[
c_0 = 0, \quad c_1 = 1, \quad c_2 = \frac{2}{2!} = 1, \quad c_3 = 0.
\]

Therefore:
\[
n^2 = 0 + n^{\underline{1}} + n^{\underline{2}}
= n + n(n-1)
= n^2 - n + n
= n^2. \checkmark
\]

\begin{remark*}[Practical value]
Once a function is expressed in the falling factorial basis, both
differencing and antidifferencing become purely mechanical, following the
rule \(\Delta\,x^{\underline{k}} = k\,x^{\underline{k-1}}\) and its
inverse. This makes the Fundamental Theorem of Finite Calculus
immediately applicable to compute sums like \(\sum_{k=0}^{n-1} k^2\)
without clever tricks.
\end{remark*}
